\documentclass[12pt,a4paper]{article}
\usepackage{ctex}
\usepackage{amsmath,amssymb,amsfonts}
\usepackage{graphicx}
\usepackage{booktabs}
\usepackage{geometry}
\geometry{left=2.5cm,right=2.5cm,top=2.5cm,bottom=2.5cm}
\usepackage{hyperref}

\title{C题 --- 微电网购电优化问题}
\date{}

\begin{document}
\maketitle

% ============================================================
\section{问题二：考虑预测误差的购电优化模型}
% ============================================================

\subsection{问题分析}

问题二要求在第一问的基础上，考虑实际场景中负载和光伏出力存在预测误差的情况。
由于日前购电计划需要提前一天制定，而实际负载和光伏发电功率在计划执行时可能与预测值存在偏差。
当实际净负荷超过计划购电量与储能调度量之和时，必须通过紧急购电来弥补缺口，
而紧急购电的电价为正常电价的5倍，这将显著增加总购电成本。

因此，本问的核心在于：
\begin{enumerate}
    \item 如何准确预测未来一天的负载和光伏发电功率；
    \item 在预测值的基础上，如何制定日前最优购电计划；
    \item 如何量化预测误差导致的紧急购电量及其成本。
\end{enumerate}

\subsection{模型假设}

\begin{itemize}
    \item 日前购电计划基于预测的负载和光伏功率制定；
    \item 实际执行时计划购电量不变，储能可在物理约束内实时充放电以平抑预测误差；
    \item 紧急购电仅在储能调节后仍有能量缺口时发生；
    \item 紧急购电电价为正常电价的5倍；
    \item 紧急购电不经过储能系统，直接弥补能量缺口；
    \item 储能系统的充放电效率、容量等参数与第一问相同。
\end{itemize}

\subsection{模型建立}

\subsubsection{两阶段优化框架}

本问题采用两阶段优化模型：

\textbf{第一阶段：日前优化（Day-ahead Optimization）}

基于预测的负载功率 $\hat{L}_t$ 和光伏功率 $\hat{P}_t$，制定次日的购电计划。
优化目标为最小化正常购电成本：

\begin{equation}
    \min_{g_t, c_t, d_t} \sum_{t=1}^{T} \left(p_t \cdot g_t + \varepsilon \cdot (c_t + d_t)\right)
\end{equation}

其中 $T = 144$ 为每天的时段数（每10分钟一个时段），$p_t$ 为时段 $t$ 的电价（元/kWh），
$g_t$、$c_t$、$d_t$ 分别为时段 $t$ 的购电量、充电量和放电量（单位均为 kWh），
$\varepsilon = 10^{-7}$ 为极小的退化惩罚系数，用于避免求解器同时分配充电和放电量。

\textbf{约束条件：}

\textbf{(1) 能量平衡约束：}
\begin{equation}
    g_t + \hat{P}_t \cdot \Delta t + d_t \geq \hat{L}_t \cdot \Delta t + c_t, \quad \forall t
\end{equation}

\textbf{(2) SOC动态约束：}
\begin{equation}
    S_{t+1} = S_t + \eta \cdot c_t - \frac{d_t}{\eta}, \quad \forall t
\end{equation}
其中 $\eta = 0.90$ 为储能系统充放电效率。

\textbf{(3) 初始SOC约束：}

题目仅规定2025年1月1日0:00的储电量为6000 kWh，此后储电量只需保持在
$[S_{\min}, S_{\max}]$ 区间内。因此储能状态在时间上是\textbf{跨日连续}的：
每天0:00的SOC等于前一天24:00的SOC，不要求每天首尾SOC相同，只对全年第一个时段施加固定初值约束：
\begin{equation}
    S_0^{(1)} = S_{\text{init}} = 6000 \text{ kWh}, \qquad
    S_0^{(d)} = S_{T}^{(d-1)}, \quad d = 2, 3, \ldots
\end{equation}
其中 $S_0^{(d)}$ 表示第 $d$ 天0:00的SOC，$S_T^{(d-1)}$ 表示前一天24:00的SOC。
该口径使储能可以在多日之间转移电量（例如低价时段多充、高价时段多放），
与题目“后继保持1200--10800 kWh”的表述一致，也带来了可观的成本下降空间。

\textbf{(4) 变量边界约束：}
\begin{equation}
    \begin{cases}
        0 \leq g_t & \text{（购电非负）} \\
        0 \leq c_t \leq P_{\max} \cdot \Delta t & \text{（充电能量上限）} \\
        0 \leq d_t \leq P_{\max} \cdot \Delta t & \text{（放电能量上限）} \\
        S_{\min} \leq S_t \leq S_{\max} & \text{（SOC安全范围）}
    \end{cases}
\end{equation}
其中 $P_{\max} = 5000$ kW，$S_{\min} = 1200$ kWh，$S_{\max} = 10800$ kWh。

\textbf{第二阶段：实时储能调度与紧急购电（Real-Time Dispatch）}

第一阶段给出的是\textbf{计划}购电量 $g_t^*$，按题目规定，实际结算时非紧急时段的
购电费用一律按计划购电量计算，因此 $g_t^*$ 一旦确定就不再变动。但储能的充放电
在实时运行中是\textbf{可调}的：当实际净负荷高于预测时，应优先由储能放电补足缺口，
而不是直接触发5倍电价的紧急购电；当实际净负荷低于预测时，多余的电力可以充入储能。
因此在每个时段按如下因果规则调度储能（只依据截至当前时段的实际值，不使用未来信息）：

\begin{equation}
    \Delta_t = \left(L_t - P_t\right) \cdot \Delta t - g_t^*, \qquad
    \left(c_t^{\text{rt}}, d_t^{\text{rt}}\right)
    = \begin{cases}
        (0, \min(\Delta_t,\, P_{\max}\Delta t,\,
                 \eta\,(S_t - S_{\min}))) & \Delta_t > 0 \quad \text{（缺电，放电）} \\[4pt]
        (\min(-\Delta_t,\, P_{\max}\Delta t,\,
               (S_{\max} - S_t)/\eta),\; 0) & \Delta_t < 0 \quad \text{（余电，充电）}
    \end{cases}
\end{equation}

\begin{equation}
    S_{t+1} = S_t + \eta\, c_t^{\text{rt}} - d_t^{\text{rt}} / \eta, \qquad
    S_{\min} \leq S_{t+1} \leq S_{\max}
\end{equation}

只有在储能充满/放空仍无法平衡时，才产生紧急购电：
\begin{equation}
    e_t = \max\left(0,\; \left(L_t - P_t\right)\Delta t + c_t^{\text{rt}}
                     - d_t^{\text{rt}} - g_t^*\right)
\end{equation}

紧急购电成本与总购电成本为：
\begin{equation}
    C_{\text{emergency}} = \sum_{t=1}^{T} 5 \cdot p_t \cdot e_t, \qquad
    C_{\text{total}} = \underbrace{\sum_{t=1}^{T} p_t \cdot g_t^*}_{\text{正常购电成本}} + C_{\text{emergency}}
\end{equation}

相比“充放电量保持计划值不变”的保守做法，实时调度把储能的调节能力真正用于
吸收预测误差，可显著削减高价紧急购电量。

\subsubsection{日前计划的安全边际：新闻童分位}

实时调度只能部分吸收误差，且放电量受剩余电量限制。由于欠购的边际成本
（5倍电价）远高于多购的边际成本（1倍电价，多余电量可存入储能），
日前计划不应取净负荷的均值，而应取\textbf{高分位数}。

设净负荷预测误差为 $\xi_t = N_t - \hat{N}_t$，欠购单位成本 $C_u = 4p_t$
（紧急价与正常价之差），多购单位成本 $C_o = p_t$，则经典报童模型给出的最优分位数为
\begin{equation}
    q^* = \frac{C_u}{C_u + C_o} = \frac{4p_t}{5p_t} = \frac{5}{6}
\end{equation}

考虑到储能本身能够免费吸收一部分缺口，实际最优分位数低于 $5/6$。
本文用预测模型自身过去 $R=30$ 天的\textbf{滚动样本外残差}经验分位数作为安全边际：
\begin{equation}
    \hat{N}_{d^*,t}^{\text{plan}} = \hat{N}_{d^*,t}
        + Q_q\left(\left\{\xi_{d,t}\right\}_{d \in \mathcal{R}}\right)
\end{equation}
其中 $\mathcal{R}$ 为目标日之前最近的 $R$ 天（严格不含目标日，无前瞻偏差），
$Q_q(\cdot)$ 为按分位数 $q$ 取的经验分位数；本文取 $q = 0.75$。

此外，日前规划时对SOC施加比物理边界更窄的\textbf{规划带}
$[S_{\min}, S_{\max}^{\text{plan}}] = [1200, 10400]$ kWh，并要求日末
$S_T \geq S_{\text{end}} = 2500$ kWh，以预留储能余量应对次日的预测误差；
实时调度仍使用物理边界 $[1200, 10800]$ kWh。

\subsubsection{基线预测方法}

为验证傅里叶谐波回归方法的优越性，本文同时实现了三种基线预测方法作为对比。

\textbf{方法一：上周同日预测法（LastWeek）}

直接使用7天前（上周同一天）的实际数据作为预测值：

\begin{equation}
    \hat{L}_{d^*,t} = L_{d^*-7,t}, \quad \hat{P}_{d^*,t} = P_{d^*-7,t}
\end{equation}

其中 $d^*$ 为目标日，$t$ 为时段索引。该方法假设一周内的用电模式具有周期性，
简单直接但无法捕获长期趋势变化。

\textbf{方法二：前4周同星期加权平均法（WeightedAverage）}

取过去4周同一星期几的数据，按时间远近赋予不同权重进行加权平均：

\begin{equation}
    \hat{L}_{d^*,t} = \frac{\sum_{k=1}^{4} w_k \cdot L_{d_k,t}}{\sum_{k=1}^{4} w_k}, \quad
    \hat{P}_{d^*,t} = \frac{\sum_{k=1}^{4} w_k \cdot P_{d_k,t}}{\sum_{k=1}^{4} w_k}
\end{equation}

其中 $d_1, d_2, d_3, d_4$ 分别为距目标日最近的4个同星期几的历史日（从近到远），
权重 $w_1=4, w_2=3, w_3=2, w_4=1$，即越近的数据权重越大。
该方法通过加权平滑降低了单一样本的随机波动，但仍局限于同一星期几的数据。

\textbf{方法三：前4周同星期线性拟合法（LinearFit）}

对过去4周同一星期几的数据，在每个时段 $t$ 上分别拟合线性趋势并外推：

\begin{equation}
    \hat{L}_{d^*,t} = a_{L,t} \cdot x^* + b_{L,t}, \quad
    \hat{P}_{d^*,t} = a_{P,t} \cdot x^* + b_{P,t}
\end{equation}

其中 $x = 1, 2, 3, 4$ 表示4个历史同星期日的序号（从远到近），$x^* = 5$ 为目标日序号。
回归系数由最小二乘法求得：

\begin{equation}
    a_{L,t} = \frac{\sum_{i=1}^{4}(x_i - \bar{x})(L_{d_i,t} - \bar{L}_t)}{\sum_{i=1}^{4}(x_i - \bar{x})^2}, \quad
    b_{L,t} = \bar{L}_t - a_{L,t} \cdot \bar{x}
\end{equation}

其中 $\bar{x} = 2.5$，$\bar{L}_t = \frac{1}{4}\sum_{i=1}^{4} L_{d_i,t}$。
该方法能捕获短期线性趋势，但对非线性变化的适应性有限。

\subsubsection{傅里叶谐波回归预测模型}

本文提出基于傅里叶谐波回归的预测方法，综合利用星期效应和季节性趋势。

\textbf{目标变量：}

定义净负荷 $N_t = L_t - P_t$，日前优化仅依赖净负荷的预测值，因此直接对净负荷进行预测。

\textbf{特征设计：}

对于第 $d$ 天，构建12维特征向量 $\mathbf{x}_d$：

\begin{equation}
    \mathbf{x}_d = \left[1,\; e_0, e_1, \ldots, e_6,\; \sin\frac{2\pi d}{365},\; \cos\frac{2\pi d}{365},\; \sin\frac{4\pi d}{365},\; \cos\frac{4\pi d}{365}\right]^T
\end{equation}

各特征含义如下：
\begin{itemize}
    \item $1$：截距项；
    \item $e_0, \ldots, e_6$：星期独热编码（$e_0$=周一，\ldots，$e_6$=周日），恰好有一个为1；
    \item $\sin(2\pi d/365), \cos(2\pi d/365)$：年度谐波，捕获季节性变化；
    \item $\sin(4\pi d/365), \cos(4\pi d/365)$：半年谐波，捕获更精细的季节性波动。
\end{itemize}

\textbf{分时段独立回归：}

一天包含 $T = 144$ 个10分钟时段。对每个时段 $t$，使用所有历史天数据独立进行最小二乘回归：

\begin{equation}
    \boldsymbol{\beta}_t = \arg\min_{\boldsymbol{\beta}} \sum_{d \in \mathcal{D}} \left(\mathbf{x}_d^T \boldsymbol{\beta} - N_{d,t}\right)^2
\end{equation}

其中 $\mathcal{D}$ 为目标日之前的所有历史天，$N_{d,t}$ 为第 $d$ 天时段 $t$ 的实际净负荷。

由于所有时段共享相同的设计矩阵 $\mathbf{X}$（特征仅依赖天，不依赖时段），
144个回归可通过一次矩阵运算同时求解：

\begin{equation}
    \mathbf{B} = (\mathbf{X}^T \mathbf{X})^{-1} \mathbf{X}^T \mathbf{Y}
\end{equation}

其中 $\mathbf{X} \in \mathbb{R}^{|\mathcal{D}| \times 12}$，$\mathbf{Y} \in \mathbb{R}^{|\mathcal{D}| \times 144}$，
$\mathbf{B} \in \mathbb{R}^{12 \times 144}$ 为系数矩阵。

\textbf{预测：}

目标日 $d^*$ 的净负荷预测值为：

\begin{equation}
    \hat{N}_{d^*,t} = \mathbf{x}_{d^*}^T \boldsymbol{\beta}_t, \quad t = 0, 1, \ldots, 143
\end{equation}

\textbf{净负荷滞后特征增强：}

上述12维特征只含天级别信息，无法刻画净负荷的短期自相关。为此在每个时段 $t$ 的回归中
额外引入前 $K=3$ 天\emph{同一时段}的实际净负荷作为滞后特征：

\begin{equation}
    \hat{N}_{d^*,t} = \mathbf{x}_{d^*}^T \boldsymbol{\beta}_t
        + \sum_{k=1}^{K} \gamma_{k,t} \cdot N_{d^*-k,\,t}
\end{equation}

滞后项的加入使设计矩阵增广，同时求解 $144$ 组最小二乘回归即可得到
$\boldsymbol{\beta}_t$ 与 $\gamma_{k,t}$。该增强把净负荷的日间惯性纳入模型，
在不引入前瞻偏差的前提下明显降低了预测误差（见下文MAE对比）。

\textbf{该方法的优势在于：}
\begin{itemize}
    \item 星期独热编码直接区分工作日与周末的用电模式差异；
    \item 年度和半年谐波捕获季节性趋势（如夏季空调负荷增加）；
    \item 144个时段独立回归，每个时段有独立的系数，天然能描述日内变化形态；
    \item 仅使用目标日之前的数据，无前瞻偏差。
\end{itemize}

\textbf{消融对比：日内谐波增强（Fourier+Intraday）。}

为检验“逐时段独立回归”的必要性，本文另构造一个消融模型：保留全部12维天级别特征，
同时显式加入日内周期（周期为144个时段）的谐波项，并把144个时段的样本展开为
一次全局最小二乘回归：

\begin{equation}
    \hat{N}_{d,t} = \mathbf{x}_d^T \boldsymbol{\beta}
        + \sum_{k=1}^{K}\left[a_k \sin\frac{2\pi k t}{144} + b_k \cos\frac{2\pi k t}{144}\right],
        \quad K = 6
\end{equation}

该模型的特征数增至 $12 + 2K = 24$ 个。由于所有时段共用同一组系数，
模型对日内形态的表达能力被大幅压缩，实验结果（见下文）证实其预测精度
显著劣于逐时段独立回归，从而反证了分时段回归设计的必要性。

\subsection{模型求解}

\subsubsection{求解算法}

\textbf{第一阶段（日前优化）：} 采用线性规划（LP）求解，使用SciPy的HiGHS求解器。
决策变量共 $4T + 1 = 577$ 个（$g_t, c_t, d_t$ 各144个，$S_t$ 共145个），
约束条件包括144条能量平衡不等式和145条SOC相关等式（144条动态等式加1条初始SOC等式）。

\textbf{第二阶段（实时调度与紧急购电）：} 将日前计划 $g_t^*$ 代入实际数据，
按式(9)--(11)逐时段执行因果储能调度，残余缺口按 $5p_t$ 结算，无需再求解优化问题。

\subsubsection{求解流程}

\begin{enumerate}
    \item 加载2025年1月1日起的历史数据（365天 $\times$ 144时段）；
    \item 对2025年2月1日至12月31日共334天，逐日进行滚动预测和优化：
    \begin{enumerate}
        \item 使用目标日之前的所有历史数据训练傅里叶+滞后预测模型；
        \item 预测目标日的净负荷曲线 $\hat{N}_t$，并叠加新闻童分位安全边际；
        \item 求解日前线性规划，得到计划购电量 $g_t^*$（规划带 $[1200, 10400]$ kWh，
              日末要求 $S_T \geq 2500$ kWh）；
        \item 实际数据到达后执行实时储能调度，得到紧急购电量 $e_t$；
        \item 累计正常购电成本和紧急购电成本，并把当日24:00的SOC传递给次日作为初值。
    \end{enumerate}
    \item 汇总全年结果。
\end{enumerate}

\subsection{模型结果}

\subsubsection{预测精度分析}

将本文提出的\textbf{傅里叶+滞后}模型与其他五种预测方法进行对比，各方法的预测精度如下表所示。
表中所有方法均只使用\textbf{点预测值}（即不含新闻童安全边际），以保证误差指标的可比性。

\begin{table}[htbp]
\centering
\caption{不同预测方法的预测精度对比（净负荷 $N=L-P$）}
\begin{tabular}{lcccc}
\toprule
\textbf{预测方法} & \textbf{MAE (kW)} & \textbf{RMSE (kW)} & \textbf{MAPE} & \textbf{特征数} \\
\midrule
上周同日（LastWeek）            & 327.77 & 486.47 & 92.58\%  & --- \\
前4周加权平均（4:3:2:1）        & 337.19 & 482.49 & 83.74\%  & --- \\
前4周线性拟合                   & 381.21 & 563.49 & 112.71\% & --- \\
傅里叶谐波回归（12特征）        & 266.62 & 385.36 & 71.17\%  & 12 \\
傅里叶+日内谐波（K=6）          & 681.40 & 839.25 & 108.23\% & 24 \\
\textbf{傅里叶+滞后（12+3）}    & \textbf{246.57} & \textbf{368.58} & \textbf{71.72\%} & 15 \\
\bottomrule
\end{tabular}
\end{table}

傅里叶+滞后模型的MAE最低，为246.57 kW，相比次优的傅里叶谐波回归（266.62 kW）
进一步降低7.5\%，相比上周同日方法（327.77 kW）降低24.8\%。
值得注意的是，直接加入日内谐波项（Fourier+Intraday）反而使MAE升至681.40 kW：
该做法把原本相互独立的144个时段回归合并为一次全局回归，弱化了模型对日内逐时段
形态的分辨能力，说明\textbf{逐时段独立回归}是傅里叶框架不可或缺的组成部分；
而滞后项的收益来自其对相邻日净负荷相关性的直接利用，因此本文最终选择傅里叶+滞后模型。

\subsubsection{购电成本分析}

各预测方法对应的全年（2025年2月1日--12月31日，共334天）购电成本如下表所示。
前六行使用各方法的点预测值制定日前计划；最后一行是\textbf{本文最终方案}——
傅里叶+滞后点预测叠加新闻童分位安全边际（$q=0.75$）：

\begin{table}[htbp]
\centering
\caption{不同预测方法对应的全年购电成本对比}
\begin{tabular}{lcccc}
\toprule
\textbf{预测方法} & \textbf{计划购电量} & \textbf{紧急购电量} & \textbf{总费用} & \textbf{紧急占比} \\
                  & \textbf{(万kWh)}    & \textbf{(万kWh)}    & \textbf{(万元)} & \textbf{(\%)}    \\
\midrule
上周同日               & 2007.50 & 58.66  & 1582.94 & 22.3\% \\
前4周加权平均          & 2002.49 & 75.95  & 1678.96 & 27.0\% \\
前4周线性拟合          & 2009.33 & 68.27  & 1642.29 & 24.9\% \\
傅里叶谐波回归         & 2012.07 & 45.74  & 1510.58 & 18.3\% \\
傅里叶+日内谐波        & 1960.84 & 173.07 & 2068.21 & 43.2\% \\
傅里叶+滞后            & 2008.67 & 38.90  & 1468.94 & 16.2\% \\
\textbf{傅里叶+滞后+新闻童分位（最终方案）}
                       & \textbf{2123.31} & \textbf{7.65} & \textbf{1357.88} & \textbf{3.5\%} \\
\bottomrule
\end{tabular}
\end{table}

在纯点预测方案中，傅里叶+滞后模型的总费用最低，为1468.94万元，相比上周同日方法
节省约114.00万元，其紧急购电占比（16.2\%）也最低，说明更准确的预测确实能减少
紧急购电的发生。但与“预测更准”的收益相比，\textbf{决策端的新闻童安全边际收益更大}：
最终方案在点预测基础上叠加分位数边际后，计划购电量由2008.67万kWh提高到2123.31万kWh
（日均多购约0.34万kWh），而紧急购电量从38.90万kWh骤降至7.65万kWh，紧急购电费用
由约238万元降至47.40万元，最终总费用\textbf{1357.88万元}，比最优的纯点预测方案
再节省111.06万元，比上周同日方法节省225.06万元。这印证了在5倍惩罚的非对称成本结构下，
\textbf{适度“多买”以对冲缺电风险}比单纯追求点预测精度更有效。

\subsection{模型评价}

\subsubsection{优点}
\begin{itemize}
    \item 两阶段优化框架清晰分离了预测不确定性和优化决策，便于分析和改进；
    \item 傅里叶+滞后模型在逐时段独立回归的基础上引入净负荷滞后特征，
          充分利用了数据的周期性、季节性和相邻日相关性，预测精度最高（MAE 246.57 kW）；
    \item 日前计划叠加新闻童分位安全边际，正确处理了紧急购电电价为5倍的
          非对称成本结构，全年总费用降至1357.88万元；
    \item 线性规划求解高效可靠，334天的全年优化可在数秒内完成；
    \item 紧急购电机制保证了供电可靠性，即使预测存在偏差也能满足负荷需求。
\end{itemize}

\subsubsection{缺点}
\begin{itemize}
    \item 紧急购电成本高昂（5倍电价），对预测误差非常敏感；
    \item 傅里叶模型假设净负荷的日内形态可由天级别特征线性描述，
          对于极端天气等异常情况的预测能力有限；
    \item 未考虑储能系统的老化成本和循环寿命约束。
\end{itemize}

% ============================================================
\section{问题四：考虑实时波动电价的购电优化模型}
% ============================================================

\subsection{问题分析}

问题四指出外网电价同样是实时波动的，要求针对波动电价建立微网当天购电策略的数学模型，
并在波动电价下重新计算问题二与问题三，分别把完整购电策略保存到
\texttt{result4-2.xlsx}（对应问题二）和 \texttt{result4-3.xlsx}（对应问题三）中。

与问题二、三相比，本问新增了\textbf{价格侧的不确定性}：电价不再是每天重复的同一条曲线，
而是随日期与时刻变化的二维价格矩阵 $p_{d,t}$。这使得购电决策除了要应对负载和光伏的
预测误差外，还要应对电价的波动与预测误差。因此本问的核心在于：

\begin{enumerate}
    \item 如何刻画“计划阶段用预测价、结算阶段用实际价”的双重价格口径；
    \item 在波动电价下，如何把问题二的两阶段优化模型推广为价格感知的日前计划；
    \item 在波动电价下，如何把问题三的滚动随机 MPC 推广为同时滚动预测“净负荷 + 电价”的决策框架。
\end{enumerate}

\subsection{模型假设}

\begin{itemize}
    \item 实际电价由附件 4 给出，是逐日、逐时段波动的确定性序列；
    \item 日前计划与 MPC 目标函数中使用的是对电价的\textbf{预测值} $\hat{p}_t$，
          而所有费用结算（计划购电费、违约金、调增购电费、紧急购电费）一律按\textbf{实际电价} $p_t^{\text{act}}$ 计算；
    \item 紧急购电电价仍为交易时刻实际电价的 5 倍；
    \item 问题三版本的二段式购电规则（调减按原价 50\% 计违约金、调增部分按 1.5 倍计）与问题三完全一致，只是其中的“原价”取实际电价；
    \item 储能参数、负载与光伏的预测误差模型、跨日 SOC 连续等假设与问题二、三相同。
\end{itemize}

\subsection{模型建立}

\subsubsection{波动电价的刻画与结算口径}

记附件 4 第 $d$ 天第 $t$ 个时段（$t=1,\ldots,T$，$T=144$）的实际电价为 $p_{d,t}^{\text{act}}$
（元/kWh）。全年统计显示其\textbf{均值与附件 1 的典型日电价几乎相同}（0.7662 元/kWh），
但\textbf{波动区间显著更宽}：附件 4 的价格落在 $[0.0076,\,1.7936]$ 元/kWh，
而附件 1 仅为 $[0.3713,\,1.3952]$ 元/kWh。这意味着低价时段更便宜、高价时段更贵，
储能的“低充高放”套利空间被放大，但预测失误的代价也随之放大。

据此，本问对每个时段区分两个价格：

\begin{equation}
    \hat{p}_t = \text{预测电价（制定计划用）}, \qquad
    p_t^{\text{act}} = \text{实际电价（费用结算用）}.
\end{equation}

实际电价与固定电价的预测方式一致，采用与问题二净负荷预测同构的\textbf{傅里叶谐波回归}模型：
以“截距 + 7 个星期独热编码 + 年/半年谐波”共 12 个特征构造设计矩阵，
对 144 个时段分别做独立的最小二乘回归，且\textbf{只使用目标日之前的历史电价}，
从而保证滚动预测的因果性。

\subsubsection{问题四--问题二版本：价格感知的两阶段优化}

\textbf{第一阶段（日前优化）：} 目标函数把问题二中的固定电价替换为预测电价 $\hat{p}_t$：

\begin{equation}
    \min_{g_t, c_t, d_t} \sum_{t=1}^{T} \left(\hat{p}_t \cdot g_t + \varepsilon \cdot (c_t + d_t)\right),
\end{equation}

约束条件（能量平衡、SOC 动态、跨日 SOC 连续、变量边界）以及新闻童分位安全边际
（$q=0.75$）、规划带 $[1200,\,10400]$ kWh、日末 $S_T \geq 2500$ kWh 的口径均与问题二完全一致。
由于目标函数首次出现价格差异，LP 会自发地把购电量向预测低价时段倾斜，
在满足能量平衡的前提下降低购电成本。

\textbf{第二阶段（实时结算）：} 计划购电量 $g_t^*$ 与实时储能调度量确定后，
按\textbf{实际电价}结算，全天费用为

\begin{equation}
    C_d = \sum_{t=1}^{T} p_{d,t}^{\text{act}} \cdot g_t^{*}
        + 5 \sum_{t=1}^{T} p_{d,t}^{\text{act}} \cdot e_t,
\end{equation}

其中 $e_t$ 为实时储能调节后仍无法补足的缺口，即紧急购电量。

\subsubsection{问题四--问题三版本：波动电价下的滚动随机 MPC}

该版本沿用问题三的滚动随机 MPC 框架：在每天 0:00、6:00、12:00、18:00 四个决策节点上，
基于情景抽样（SAA）求解两阶段随机规划，仅执行未来 6 小时（36 个时段）的一阶段决策，
随后到达下一节点再重新优化，SOC 跨日连续。区别在于目标函数与结算均改用波动电价：

\textbf{节点优化目标（用预测电价）：}
\begin{equation}
    \min \; \sum_{t}
    \Big[\, \hat{p}_t \, g_t^{a} \;+\; 0.5\,\hat{p}_t \, (u_t + v_t) \;+\; 5\,\hat{p}_t \, e_t \,\Big]
    \;+\; \lambda \, \mathrm{CVaR}_{\alpha},
\end{equation}
其中 $g_t^{a}$ 为最终执行的购电量，$u_t, v_t \geq 0$ 为相对计划量 $g_t^{0}$ 的正、负偏差
（用于线性化调减违约金与调增加价），$e_t$ 为缺口紧急购电量，$\lambda$ 为风险权重（本问取 0）。

\textbf{实际结算（用实际电价）：}
\begin{equation}
    C_d = \sum_{t=1}^{T}
    \Big[\, p_{d,t}^{\text{act}} \, g_t^{a}
        \;+\; 0.5\, p_{d,t}^{\text{act}} \,\big| g_t^{0} - g_t^{a} \big|
        \;+\; 5\, p_{d,t}^{\text{act}} \, e_t \,\Big].
\end{equation}

上面两个式子结构完全对称，只是把计划价 $\hat{p}_t$ 换成了实际价 $p^{\text{act}}_{d,t}$，
从而在数学上完整实现了“\textbf{预测价做计划、实际价做结算}”的波动电价策略。

\subsection{模型求解}

\subsubsection{求解算法}

两个版本均归结为线性规划，统一使用 SciPy 的 HiGHS 求解器：

\begin{itemize}
    \item \textbf{问题二版本}：单日一个 LP，决策变量约 ${4T+1}$ 个，一次性求解全日计划；
    \item \textbf{问题三版本}：每个决策节点求解一个 SAA 两阶段 LP
          （一阶段 $g^0, c, d, S$ 共 $4\times 36$ 个非预期性变量，
           二阶段按 12 个情景展开 $g^a, c, d, S, u, v, e$），
          全天共 4 个节点、约 30 次求解，334 天合计约 1.2 万次 LP。
\end{itemize}

\subsubsection{求解流程}

\begin{enumerate}
    \item 加载附件 4 的实际电价矩阵 $p^{\text{act}}_{d,t}$（365 天 $\times$ 144 时段），
          并用“目标日之前”的历史数据训练傅里叶电价预测模型；
    \item 对 2025 年 2 月 1 日至 12 月 31 日共 334 天逐日滚动：
    \begin{enumerate}
        \item 预测目标日的电价 $\hat{p}_t$（以及净负荷、光伏）；
        \item 问题二版本：求解一次日前 LP 得到 $g_t^*$，再用实际电价结算；
              问题三版本：在四个节点上依次求解随机 MPC 并因果执行储能；
        \item 把当日 24:00 的 SOC 传递给次日作为初值；
    \end{enumerate}
    \item 汇总全年结果，分别写入 \texttt{result4-2.xlsx}（计划购电量/充放电量/紧急购电量）
          与 \texttt{result4-3.xlsx}（计划购电量/调整购电量/充放电量/紧急购电量）。
\end{enumerate}

\subsection{模型结果}

\subsubsection{波动电价下的全年购电结果}

\begin{table}[htbp]
\centering
\small
\caption{问题四两个版本的全年购电结果（2025-02-01--12-31，共334天）}
\begin{tabular}{lcccc}
\toprule
\textbf{方案} & \textbf{计划购电量} & \textbf{紧急购电量} & \textbf{总费用} & \textbf{对应文件} \\
              & \textbf{(万kWh)}    & \textbf{(万kWh)}    & \textbf{(万元)} & \\
\midrule
问题二版本（两阶段 LP）    & 2123.19 & 7.67  & 1431.34 & result4-2.xlsx \\
问题三版本（滚动随机 MPC） & 2019.73 & 10.43 & 1501.35 & result4-3.xlsx \\
\bottomrule
\end{tabular}
\end{table}

其中问题三版本在第一阶段计划之外，还会在实时运行中把购电量调整到 2141.10 万kWh
（调整幅度约 121.37 万kWh）。问题二版本的 1431.34 万元可分解为正常购电费 1381.01 万元
与紧急购电费 50.33 万元；问题三版本的 1501.35 万元可分解为计划购电费 1298.46 万元、
调减违约金 $-3.03$ 万元、调增购电费 140.23 万元与紧急购电费 65.69 万元。

\subsubsection{与固定电价情形的对比}

\begin{table}[htbp]
\centering
\small
\caption{固定电价（问题二、三）与波动电价（问题四）的全年总费用对比}
\begin{tabular}{lccc}
\toprule
\textbf{方案} & \textbf{固定电价} & \textbf{波动电价} & \textbf{增幅} \\
              & \textbf{(万元)} & \textbf{(万元)} & \\
\midrule
问题二版本（两阶段 LP）    & 1357.88 & 1431.34 & $+73.46$（$5.4\%$） \\
问题三版本（滚动随机 MPC） & 1427.33 & 1501.35 & $+74.02$（$5.2\%$） \\
\bottomrule
\end{tabular}
\end{table}

由表可见：

\begin{enumerate}
    \item \textbf{波动电价使总费用上升约 5\%}（约 73--74 万元）。由于附件 4 与附件 1 的
          价格均值几乎相同（均为 0.7662 元/kWh），这一增幅并非来自平均价格水平，
          而是来自\textbf{价格波动本身}：峰值电价更高（1.7936 元/kWh 对 1.3952 元/kWh），
          而储能受容量与功率上限约束，无法完全把购电量搬运到谷价时段，
          残余的高价时段购电以及 5 倍紧急电价带来的费用上升共同推高了总成本。
    \item \textbf{计划购电量几乎不变}（问题二版本由 2123.31 万kWh 变为 2123.19 万kWh）。
          这是因为能量平衡约束决定了全年的总购电量由总负荷、总光伏和储能损耗决定，
          与价格无关；价格波动只改变购电量在时段间的\textbf{分布}，不改变其\textbf{总量}。
    \item \textbf{紧急购电量保持低位}（问题二版本 7.67 万kWh），其对应的紧急购电费
          （50.33 万元）仅占总费用的 3.5\%。说明新闻童分位安全边际在波动电价下依然有效，
          缺电风险仍被有效对冲。
    \item 问题三版本的费用高于问题二版本（1501.35 万元对 1431.34 万元），
          与固定电价下的结论一致：滚动随机 MPC 需要为 6 小时内的不确定性预留裕度，
          并承担“调增购电 1.5 倍”的额外成本，因此其保守性带来了一定费用溢价，
          但换取了在任意时点都可重新决策、对预报更新更敏感的灵活性。
\end{enumerate}

\subsection{模型评价}

\subsubsection{优点}

\begin{itemize}
    \item 通过“预测价做计划、实际价做结算”的双价格口径，清晰地把价格不确定性
          拆分为决策与结算两个环节，模型结构对称、易于实现；
    \item 价格预测复用了净负荷预测的傅里叶谐波回归框架，工程实现统一；
    \item 两个版本分别对应问题二与问题三的建模路线，结果可直接对照，
          完整回答了“在波动电价下重新计算问题 2 和问题 3”的要求；
    \item 结果文件 \texttt{result4-2.xlsx} 与 \texttt{result4-3.xlsx} 的时段划分、
          列序与附件 5 模板严格一致。
\end{itemize}

\subsubsection{缺点}

\begin{itemize}
    \item 电价预测仅使用时间特征（星期、季节谐波），未引入负荷、温度等外部解释变量，
          对突发价格尖峰的预测能力有限；
    \item 采用点预测价进入目标函数，未把价格预测误差显式纳入随机规划，
          若电价的预测误差显著增大，计划与结算的偏离会进一步扩大；
    \item 未考虑电价与负荷之间的相关性（两者在实际系统中往往同向波动）。
\end{itemize}

\end{document}
